\chapter
[\CO{[ALS]} Erweiterte Reanimation (Advanced Life Support)]
{\CC{[ALS]}\\Erweiterte Reanimation\\(Advanced Life Support)}
\label{chp:ALS}

\minitoc

% Die Durchführung von ALS-Maßnahmen ist möglich mit Sanitätern aber der
% Qualifikation NKV oder (Not-)Ärzten. Ansonsten muß eine Basisreanimation (wenn
% vorhanden mittels AED) durchgeführt und weitere Hilfe angefordert werden. 
% 
% \section{Einleitung}
% 
% ERC
% 
% \section{Eingesetzte Materialien}
% 
% \subsection{Medikamente für die Reanimation}
% 
% \begin{table}[H]
% \begin{gWideParEnv}
% \Captionof{table}{Übersicht: Die wichtigsten bei der Reanimation eingesetzten
% Medikamente.}
% 
% \begin{tabulary}{\linewidth}{lLLL}
% \addlinespace
% \noindent\TabHeading{Wirkstoff}
% &
% \noindent\TabHeading{Erläuterungen}
% &
% \noindent\TabHeading{Indikation}
% &
% \noindent\TabHeading{Spezialitäten}
% \\\midrule\addlinespace
% \TabSub{Adrenalin}
% &
% Vasokonstriktor
% &
% Schockbarer Rhythmus
% 
% Nicht-Schockbarer Rhythmus
% &
% L-Adrenalin Fresenius 2\,mg / 20\,ml
% 
% L-Adrenalin Fresenius 0,5\.mg / 5\,ml
% 
% Suprarenin 1\,mg / 1\,ml
% \\\addlinespace
% \TabSub{Atropin}
% &
% Anti\-para\-sym\-patho\-mi\-meti\-kum
% &
% Nicht-Schockbarer Rhythmus
% &
% Atropinum sulfuricum Nycomed 0,5\,mg auf 1\,ml
% \\\addlinespace
% \TabSub{Amiodaron}
% &
% Anti\-ar\-rhythmi\-kum 
% 
% Kl. III
% &
% Schockbarer Rhythmus
% &
% Sedacoron 150\,mg auf 3\,ml
% \\\addlinespace\bottomrule
% \end{tabulary}
% 
% \end{gWideParEnv}
% \end{table}
% 
% \subsubsection{Adrenalin}
% 
% \Parallel{
% L-Adrenalin Fresenius 2\,mg / 20\,ml
% 
% L-Adrenalin Fresenius 0,5\.mg / 5\,ml
% 
% Suprarenin 1\,mg / 1\,ml
% }{
% \Captionof{table}{Adrenalingabe bei der Reanimation}
% 
% \begin{tabulary}{\linewidth}{LL}
% \TabHeading{Indikation}
% &
% \TabHeading{Dosierung}
% \\\midrule\addlinespace
% \TabSub{Schockbarer Rhythmus}
% &
% 1\,mg i.v. 
% 
% vor dem 3., 5., 7.,~{\ldots} Schock
% \\\addlinespace
% \TabSub{Nicht-Schockbarer Rhythmus}
% &
% 1\,mg i.v.
% 
% sobald venöser Zugang geschaffen wurde alle 3-5 Minuten
% \\\addlinespace\bottomrule
% \end{tabulary}
% }
% 
% \subsubsection{Atropin}
% 
% \Parallel{
% 
% }{
% \Captionof{table}{Atropingabe bei der Reanimation}
% 
% \begin{tabulary}{\linewidth}{LL}
% \TabHeading{Indikation}
% &
% \TabHeading{Dosierung}
% \\\midrule\addlinespace
% \TabSub{Schockbarer Rhythmus}
% &
% keine Gabe
% \\\addlinespace
% \TabSub{Nicht-Schockbarer Rhythmus}
% &
% 1\,mg i.v.
% 
% sobald venöser Zugang geschaffen wurde alle 3-5 Minuten
% \\\addlinespace\bottomrule
% \end{tabulary}
% }
% 
% \subsubsection{Amiodaron}
% 
% \Parallel{
% 
% }{
% \Captionof{table}{Amiodarongabe bei der Reanimation}
% 
% \begin{tabulary}{\linewidth}{LL}
% \TabHeading{Indikation}
% &
% \TabHeading{Dosierung}
% \\\midrule\addlinespace
% \TabSub{Schockbarer Rhythmus}
% &
% 1\,mg i.v. 
% 
% vor dem 3., 5., 7.,~{\ldots} Schock
% \\\addlinespace
% \TabSub{Nicht-Schockbarer Rhythmus}
% &
% keine Gabe
% \\\addlinespace\bottomrule
% \end{tabulary}
% }
% 
% \subsection{Geräte}
% 
% \subsubsection{Beatmungsgeräte}
% 
% Atemzugsvolumen (AZV) / Atemminutenvolumen (AMV)
% 
% \subsubsection{Defibrillatoren}
% 
% AED-Modus
% 
% Manueller Modus
% 
% Sync-Taste
% 
% \section{Problemstellungen}
% 
% \subsection{Atemweg}
% 
% Beutel-Masken-Beatmung
% 
% Erweitertes Atemwegsmanagement
% 
% Endotracheale Intubation
% 
% \See{topic:intubation-assistenz}
% 
% \subsection{Thoraxmassage}
% 
% \section{Ablauf}
% 
% \begin{enumerate}
%     \item Basisreanimation
%     \item Wenn ALS noch nicht möglich ist, ALS-Maßnahmen wenn möglich
%     vorbereiten.
%     \item Eintreffen von ALS-Team
%     \item Erweiterte Maßnahmen
%     \begin{itemize}
%         \item Venöser Zugang
%         \item Medikamentengabe
%         \item Erweitertes Atemwegsmanagement
%         \item Sonstige Maßnahmen
%     \end{itemize}
% \end{enumerate}

\clearpage
\begin{figure}
\begin{gWideParEnv}
\Captionof{figure}{Algorithmus Advanced Live Support (ALS)}
\includegraphics[width=\linewidth]{./extern/ALS.pdf}
\end{gWideParEnv}
\end{figure}

\nocite{Uray2008}
